\documentclass[12pt]{article}

\usepackage{sbc-template}

\usepackage{graphicx,url}

\usepackage[brazil]{babel}
\usepackage[utf8]{inputenc}

\usepackage{float}
\usepackage{tikz}
\usetikzlibrary{arrows.meta,positioning}

\sloppy

\title{WhyRunner: Juiz Online Educacional com Execução Isolada\\e Feedback por Inteligência Artificial}

\author{Juan Israel dos Anjos}

\address{Instituto Federal de Educação, Ciência e Tecnologia de Minas Gerais
  (IFMG)\\
  Rua Afonso Sardinha, 90, Bairro Minas Talco\\ Ouro Branco/MG -- CEP: 36494-018 -- Brasil
  \email{juaniwk3@gmail.com}
}

\begin{document}

\maketitle

\begin{abstract}
  Most online judges were designed for contests and return only terse verdicts such as ``Wrong Answer''. This paper presents WhyRunner, an online judge for programming education that combines sandboxed execution in six languages, professor-curated classes and lessons, AI-assisted problem authoring with mandatory reference-solution validation, and on-demand hints from a large language model. An automated evaluation of the judge measured latency per language and the containment of 12 malicious-program scenarios: all languages were graded correctly, five in under 0.4\,s for a 10-test problem, and every attack was contained, while a slow Portugol runtime and an unbounded program output were found. Classroom evaluation is left as future work.
\end{abstract}

\begin{resumo}
  A maioria dos juízes online foi concebida para competições e devolve apenas veredictos sucintos, como ``Resposta Errada''. Este artigo apresenta o WhyRunner, um juiz online voltado ao ensino de programação que combina execução isolada em seis linguagens, turmas e aulas curadas pelo professor, autoria de problemas assistida por IA com validação obrigatória por solução de referência e dicas sob demanda geradas por um modelo de linguagem. Uma avaliação automatizada do juiz mediu a latência por linguagem e a contenção de 12 cenários de código malicioso: todas as linguagens foram corrigidas corretamente, cinco em menos de 0,4\,s para problemas de 10 casos, e todos os ataques foram contidos, mas revelaram-se lentidão no Portugol e ausência de limite para a saída dos programas. A avaliação em sala de aula fica como trabalho futuro.
\end{resumo}

\section{Introdução}

A programação competitiva e a resolução de listas de exercícios são práticas consolidadas no ensino de Computação, pois exercitam o raciocínio lógico, a modelagem de problemas e o domínio de algoritmos. Para viabilizar essas práticas em turmas numerosas, popularizaram-se os \emph{juízes online}: sistemas que recebem o código do estudante, executam-no contra casos de teste predefinidos e devolvem um veredicto automaticamente \cite{wasik2018survey}.

O problema é que a maioria desses sistemas nasceu no contexto de maratonas, priorizando justiça e eficiência em detrimento do aspecto pedagógico: o retorno se limita a mensagens como ``Resposta Errada'', sem indicar \emph{por que} a solução falhou \cite{keuning2018systematic}, e poucos permitem ao professor montar atividades a partir de problemas existentes. Como o \emph{feedback} é um dos fatores de maior impacto na aprendizagem \cite{hattie2007power}, e os modelos de linguagem de grande porte (LLMs) tornaram viável explicar o erro sem entregar a solução \cite{liffiton2023codehelp,kazemitabaar2024codeaid}, há espaço para um juiz voltado ao ensino.

O objetivo geral deste trabalho é desenvolver e avaliar o \textbf{WhyRunner}, um juiz online voltado ao ensino de programação. Os objetivos específicos são: (i) executar código submetido em seis linguagens de forma segura e isolada; (ii) oferecer ao professor um modelo de turmas, aulas e exercícios, competições em tempo real e ferramentas de autoria de problemas assistidas por IA; (iii) fornecer ao estudante dicas geradas por IA que orientem sem revelar a resposta; e (iv) avaliar o desempenho e a segurança do módulo de correção.

\section{Referencial teórico}\label{sec:referencial}

\subsection{Juízes online e isolamento}

Um juiz online compila e executa o código do estudante em ambiente homogêneo e o compara contra casos de teste, devolvendo veredictos como \emph{Aceito}, \emph{Resposta Errada} ou \emph{Tempo Limite Excedido} \cite{wasik2018survey}. Para que os casos ocultos não sejam extraídos um a um, sistemas de competição como o BOCA mostram apenas o veredicto \cite{de2004boca}. Como o código é arbitrário, cada submissão roda em uma \emph{sandbox}: o \texttt{isolate} \cite{mares2012isolate}, criado para a Olimpíada Internacional de Informática (IOI), confina o processo com \emph{namespaces} do Linux (rede, processos, arquivos) e limites de tempo, memória e processos, sem o \emph{daemon} privilegiado exigido por contêineres Docker \cite{merkel2014docker}.

\subsection{Feedback e IA no ensino de programação}

Um veredicto binário diz ao estudante onde ele está, mas não qual o próximo passo \cite{hattie2007power,keuning2018systematic}. O CodeHelp \cite{liffiton2023codehelp} e o CodeAid \cite{kazemitabaar2024codeaid} usam LLMs com \emph{guardrails} que impedem a entrega de código pronto, mas são assistentes independentes, sem acesso ao caso de teste que falhou. Para iniciantes, recomenda-se começar por ambientes visuais \cite{weintrop2015block} e por \emph{Parsons problems} \cite{parsons2006puzzles}, em que o estudante ordena linhas embaralhadas de uma solução.

\subsection{Trabalhos relacionados}\label{sec:relacionados}

O \textbf{BOCA} \cite{de2004boca} é o padrão das maratonas brasileiras; o \textbf{Run.codes} \cite{runcodes2024} permite a professores criar listas e competições em várias linguagens; o \textbf{Ataux} \cite{silva2021ataux} gera planilhas de listas resolvidas no Codeforces; e o \textbf{GOJ} \cite{lima2022gamma} armazena e corrige questões das maratonas da UnB. A Tabela~\ref{tab:comparacao} resume a comparação.

\begin{table}[!htb]
\centering
\caption{Comparação entre juízes online}
\label{tab:comparacao}
\footnotesize
\setlength{\tabcolsep}{4pt}
\renewcommand{\arraystretch}{1.25}
\begin{tabular}{|l|c|c|c|c|c|c|}
\hline
\rule{0pt}{5.6ex}\textbf{Sistema} & \shortstack{\textbf{Feedback}\\\textbf{por IA}} & \textbf{Linguagens} & \textbf{Competição} & \shortstack{\textbf{Aulas}\\\textbf{curadas}} & \textbf{Restrições} & \shortstack{\textbf{Portugol}\\\textbf{+ blocos}} \\
\hline
Ataux      & Não & Limitado   & Não      & Não & Não & Não \\
GOJ        & Não & Moderado   & Limitado & Não & Não & Não \\
Run.codes  & Não & Amplo      & Sim      & Não & Não & Não \\
BOCA       & Não & Moderado   & Sim      & Não & Não & Não \\
WhyRunner  & Sim & Amplo (6)  & Sim      & Sim & Sim & Sim \\
\hline
\end{tabular}
\end{table}

Nenhum dos sistemas analisados combina correção automática com dicas por IA cientes do caso de teste que falhou, nem permite ao professor montar aulas a partir de problemas existentes, exigir uma forma de resolução (restrições de solução) ou oferecer aos iniciantes respostas em blocos que geram Portugol.

\section{Metodologia}\label{sec:metodologia}

\subsection{Processo de desenvolvimento}

O sistema foi desenvolvido de forma iterativa e incremental desde abril de 2025, em ciclos de implementação, integração, teste e refinamento. Os requisitos iniciais vieram de revisão da literatura, conversas informais com estudantes e professores e da experiência do autor com plataformas de programação competitiva. Um protótipo foi usado informalmente por cinco colegas em um servidor do Discord, sem roteiro nem questionário; três reclamações viraram mudanças (código de cada tentativa visível nas competições, placar no estilo das maratonas e erros detalhados na importação de problemas), e os elogios citaram o visual e o nome, lido como o ``uai'' mineiro. Por ser pequeno e informal, esse teste não substitui a avaliação com usuários.

Durante o desenvolvimento, a fila de submissões passou por BullMQ e AWS SQS antes de dar lugar ao \texttt{LISTEN/NOTIFY} do PostgreSQL, e a execução migrou de contêineres Docker, cujo \emph{socket} equivale a acesso de superusuário, para o \texttt{isolate}.

\subsection{Arquitetura}

O WhyRunner é composto por dois serviços que compartilham um banco PostgreSQL (Figura~\ref{fig:arquitetura}): a \textbf{aplicação web}, em Next.js com Drizzle ORM, responsável pela interface, regras de negócio e chamadas aos modelos de linguagem; e o \textbf{juiz}, em Rust \cite{matsakis2014rust}, que executa as submissões no \texttt{isolate} e grava o resultado.

\begin{figure}[!htb]
\centering
\begin{tikzpicture}[
  node distance=0.6cm and 1.2cm,
  box/.style={draw, rounded corners, fill=gray!10, text width=3.2cm, align=center, minimum height=1cm, font=\footnotesize},
  arr/.style={-{Latex[length=2mm]}, thick}
]
\node[box] (user) {Estudante / Professor (navegador)};
\node[box, right=of user] (web) {Aplicação web\\(Next.js)};
\node[box, right=of web] (llm) {IA: Gemini 3.5 Flash $\rightarrow$ Flash-Lite $\rightarrow$ Claude Haiku 5.5};
\node[box, below=of web] (db) {PostgreSQL\\(\texttt{LISTEN/NOTIFY})};
\node[box, below=of db] (judge) {Juiz (Rust)};
\node[box, right=of judge] (box) {\texttt{isolate}: uma \emph{box} por caso de teste};

\draw[arr] (user) -- (web);
\draw[arr] (web) -- (llm);
\draw[arr] (web) -- node[left, font=\scriptsize]{submissão + \texttt{NOTIFY}} (db);
\draw[arr] (db) -- node[left, font=\scriptsize]{notificação} (judge);
\draw[arr] (judge) -- (box);
\draw[arr, dashed] (judge.east) ++(0,0.3) to[bend right=40] node[right, font=\scriptsize]{resultado} (db.east);
\end{tikzpicture}
\caption{Arquitetura do WhyRunner}
\label{fig:arquitetura}
\end{figure}

O fluxo de uma submissão é: (1) a aplicação web grava o código com status \texttt{PENDING} e emite um \texttt{NOTIFY}; (2) o juiz reivindica a submissão com bloqueio de linha (\texttt{FOR UPDATE SKIP LOCKED}), o que permite vários juízes em paralelo; (3) o código é compilado uma única vez e roda em uma \emph{box} nova por caso de teste; (4) o resultado é gravado e o estudante, notificado, podendo pedir uma dica: o modelo recebe o enunciado, o código e o caso que falhou, com instruções para não fornecer a solução, e o servidor oculta dos estudantes os casos de teste secretos. As chamadas de IA tentam, em ordem, o Gemini 3.5 Flash \cite{gemini2023}, o Gemini 3.5 Flash-Lite e o Claude Haiku 5.5 \cite{anthropic2026haiku}, com limite de 30\,s por tentativa. As linguagens suportadas são C, C++, Java, Python, Rust e Portugol.

\subsection{Funcionalidades}\label{sec:funcionalidades}

\textbf{Autoria de problemas.} A IA gera uma \emph{narrativa} para o enunciado, sugere \emph{tags} e revisa a descrição, sugerindo casos de borda adicionáveis com um clique; uma pré-visualização mostra o problema como o estudante o verá. Antes de publicar, o problema precisa ser \emph{validado}: uma solução de referência, colada ou escrita pela IA a partir da descrição e dos casos salvos, roda no juiz contra todos os casos, e qualquer alteração posterior nos casos exige nova validação.

\textbf{Competições e resolução.} Problemas têm códigos curtos (ex.: \texttt{DEMO-101}) e podem ser filtrados por dificuldade e \emph{tags}. Uma competição pode ser privada e tem um nível máximo de ajuda da IA (desligada, conceito, dica ou precisa), imposto pelo servidor; o placar segue o estilo das maratonas, com tentativas, minuto da resolução e penalidade. O editor tem destaque de sintaxe e salvamento automático, e cada envio guarda o código submetido, com o botão ``Explicar'' para pedir a dica.

\textbf{Turmas e aulas.} Os estudantes entram na turma por um código. Em cada aula, o professor ordena exercícios a partir de problemas existentes, com prazo, e escolhe o modo de resposta: código, blocos (Blockly, convertidos em Portugol) ou \emph{Parsons}. Também pode impor \emph{restrições de solução}, como limite de laços aninhados (por análise estática) ou um algoritmo exigido (verificado por IA), e revisa os envios após o prazo.

\subsection{Ambiente de execução}

Cada caso de teste roda em uma \emph{box} do \texttt{isolate} nova, sem rede, com no máximo 64 processos, 128\,MB de memória para C, C++ e Python (mais para Rust e linguagens da JVM) e limite de CPU derivado do limite do problema (1\,s por padrão), ajustado por linguagem: inalterado para C, C++ e Rust, dobrado mais 0,5\,s para Java, triplicado para Python e dobrado mais 3\,s para Portugol; o tempo real é o dobro do de CPU mais 1\,s. Tempos, memória e causa de término são lidos do arquivo de metadados do \texttt{isolate}.

\subsection{Protocolo de avaliação preliminar}\label{sec:protocolo}

Como o sistema ainda não foi usado em sala de aula, esta primeira avaliação se concentra no módulo de correção, por ser o componente crítico para segurança e desempenho. Foram realizados três experimentos automatizados, todos chamando a mesma função de correção usada em produção. (1) \textbf{Correção funcional}: a suíte de testes automatizados do juiz, que cobre veredictos, limites de tempo, compilação, política de parada e restrições de solução. (2) \textbf{Desempenho por linguagem}: um programa que lê e imprime um inteiro foi corrigido 10 vezes contra 1 e contra 10 casos de teste, após um aquecimento, medindo o tempo total da correção, o tempo de execução e o pico de memória. (3) \textbf{Robustez}: 12 programas maliciosos tentaram acessar a rede, ler e escrever arquivos do sistema, matar processos, criar processos indefinidamente, esgotar memória, executar indefinidamente e inundar a saída; o ataque é considerado contido se o programa for interrompido pelo juiz ou imprimir \texttt{BLOCKED}.

Os experimentos rodaram na imagem Docker de produção do juiz (\texttt{isolate} 2.7, Python 3.11, GCC 12.2, OpenJDK 21), com os testes executados sequencialmente, em um computador com processador AMD Ryzen 5 5600 (6 núcleos, 12 \emph{threads}), 32\,GB de RAM e Linux 7.2.

\section{Resultados preliminares}\label{sec:resultados}

\subsection{Correção funcional}

Todos os 29 testes automatizados do juiz passaram: 4 do leitor de metadados do \texttt{isolate}, 12 das regras de restrição de solução e 13 de ponta a ponta no \emph{sandbox} real, que confirmam a saída correta nas seis linguagens, o respeito ao limite de tempo do problema, a preservação das mensagens do compilador e a política de parada (aulas executam todos os casos; competições param na primeira falha).

\subsection{Desempenho por linguagem}

A Tabela~\ref{tab:desempenho} apresenta o tempo médio de correção de um programa trivial com 1 caso de teste ($T_1$) e com 10 casos ($T_{10}$). A coluna ``por caso extra'' estima o custo de cada caso de teste extra, calculado como $(T_{10} - T_1)/9$; a diferença entre $T_1$ e o tempo de execução reportado pelo juiz corresponde, aproximadamente, à compilação.

\begin{table}[!htb]
\centering
\caption{Tempo de correção por linguagem (média $\pm$ desvio-padrão de 10 repetições)}
\label{tab:desempenho}
\small
\begin{tabular}{|l|r|r|r|r|}
\hline
\textbf{Linguagem} & $T_1$ (ms) & $T_{10}$ (ms) & \textbf{Por caso extra} (ms) & \textbf{Memória} (MB) \\
\hline
C        & $42 \pm 1$      & $98 \pm 3$       & 6     & 1,5 \\
C++      & $270 \pm 8$     & $335 \pm 7$      & 7     & 3,2 \\
Rust     & $265 \pm 5$     & $369 \pm 17$     & 12    & 1,8 \\
Java     & $684 \pm 25$    & $1.263 \pm 22$   & 64    & 41,8 \\
Python   & $14 \pm 0$      & $148 \pm 4$      & 15    & 7,9 \\
Portugol & $1.190 \pm 28$  & $11.601 \pm 99$  & 1.157 & 85,6 \\
\hline
\end{tabular}
\end{table}

\subsection{Robustez a código malicioso}

A Tabela~\ref{tab:ataques} apresenta o resultado dos 12 cenários de ataque. Todos foram contidos: nenhum programa acessou a rede, leu ou alterou arquivos do sistema, afetou processos fora da \emph{box} ou deixou de ser interrompido.

\begin{table}[!htb]
\centering
\caption{Cenários de ataque e resultado no \emph{sandbox}}
\label{tab:ataques}
\small
\begin{tabular}{|l|c|l|r|}
\hline
\textbf{Ataque} & \textbf{Ling.} & \textbf{Resultado} & \textbf{Tempo} (ms) \\
\hline
Conexão TCP externa              & Python & Bloqueado (sem rede)                     & 40 \\
Resolução de DNS                 & Python & Bloqueado (sem rede)                     & 30 \\
Escrita em \texttt{/etc}         & Python & Bloqueado (sistema somente leitura)      & 18 \\
Leitura de \texttt{/etc/shadow}  & Python & Bloqueado (sem permissão)                & 18 \\
Matar o processo 1 (\texttt{init}) & Python & Bloqueado (\emph{namespace} próprio)   & 23 \\
\emph{Fork bomb}                 & Python & Erro de execução (limite de processos)   & 64 \\
\emph{Fork bomb}                 & C      & Tempo limite excedido                    & 3.181 \\
Laço infinito                    & C      & Tempo limite excedido                    & 3.147 \\
Espera indefinida (\texttt{sleep}) & Python & Tempo limite excedido (tempo real)     & 7.105 \\
Alocar 500\,MB                   & Python & Erro de execução (\texttt{MemoryError})  & 14 \\
Alocar 2\,GB                     & C      & Bloqueado (\texttt{malloc} falha)        & 51 \\
Inundação da saída padrão        & C      & Tempo limite excedido                    & 7.930 \\
\hline
\end{tabular}
\end{table}


\subsection{Discussões}

Para as linguagens compiladas nativas e para Python, um problema com 10 casos é corrigido em menos de 0,4\,s, tempo compatível com o uso em sala de aula. Compilar uma única vez por submissão mostrou-se acertado: em C++ e Rust, a compilação ($\approx$\,260\,ms) seria repetida a cada caso. O gargalo é o Portugol, cujo interpretador recompila o programa em toda execução: cerca de 1,2\,s por caso e 11,6\,s para 10 casos, quase 35 vezes o tempo de C++. Como o Portugol e o modo de blocos se destinam aos iniciantes, para quem a espera é mais desmotivadora, compilá-lo uma única vez é a correção prioritária.

Quanto à segurança, os \emph{namespaces} e limites do \texttt{isolate} \cite{mares2012isolate} contiveram todos os ataques, dando ao WhyRunner o isolamento usado na IOI. A exceção parcial foi a inundação da saída padrão, interrompida só após 7,9\,s, acima do limite de 3\,s: o juiz lê toda a saída para a memória, sem limite, e deve passar a interromper o programa ao ultrapassar um tamanho máximo.

Entre os trabalhos relacionados, só o Ataux \cite{silva2021ataux} relata avaliação, de usabilidade; nenhum mede desempenho ou segurança da correção, o que torna esta avaliação complementar. Os resultados vêm de um único computador, com programas triviais e sem concorrência, e não incluem a latência do banco e da interface; as dicas por IA ainda não foram avaliadas.

\section{Conclusões}\label{sec:conclusoes}

Este artigo apresentou o WhyRunner, um juiz online voltado ao ensino que integra execução isolada em seis linguagens, turmas e aulas curadas pelo professor, autoria de problemas assistida por IA com validação obrigatória e dicas por IA a partir do caso que falhou. Os resultados preliminares mostram que o módulo de correção atende aos objetivos (i) e (iv), e revelaram dois pontos a corrigir: a lentidão do Portugol e a saída sem limite.

A principal limitação é a ausência de avaliação com usuários. Os próximos passos são: (1) aplicar o sistema em uma turma, com uma competição e uma aula, questionário SUS \cite{brooke1996sus} com os estudantes e registro da visão do professor ao criar a aula e a competição, revisar os envios e dosar a IA; (2) avaliar se as dicas revelam a solução; (3) compilar o Portugol uma única vez e limitar a saída lida pelo juiz; e (4) ativar o modo de \emph{cgroups} do \texttt{isolate} e medir a vazão com submissões concorrentes.

\bibliographystyle{sbc}
\bibliography{refs}

\end{document}
